% ACL Paper Skeleton (Long Paper — 8 pages)
% Replace placeholders marked with <...> before use
\documentclass[11pt]{article}

% ACL style — update year as needed
\usepackage[hyperref]{acl2025}

\usepackage{times}
\usepackage{latexsym}
\usepackage[T1]{fontenc}
\usepackage[utf8]{inputenc}
\usepackage{microtype}
\usepackage{inconsolata}
\usepackage{booktabs}
\usepackage{amsmath}
\usepackage{amssymb}
\usepackage{graphicx}
\usepackage{xcolor}

\newcommand{\ours}{\textsc{<MethodName>}}
\newcommand{\todo}[1]{\textcolor{red}{[TODO: #1]}}

\title{<Paper Title>}

\author{
  <Author 1> \\
  <Affiliation 1> \\
  \texttt{<email1>} \\\And
  <Author 2> \\
  <Affiliation 2> \\
  \texttt{<email2>} \\
}

\begin{document}

\maketitle

\begin{abstract}
% NLP-style abstract: task, challenge, approach, results
\todo{Write abstract — 150-250 words}
\end{abstract}

%% ============================================================
\section{Introduction}
\label{sec:intro}

\todo{Motivation — why this NLP problem matters}

\todo{Specific challenge and gap in existing work}

\todo{Our approach and linguistic/computational insight}

\todo{Results summary}

Our contributions include:
\begin{itemize}
    \item \todo{Contribution 1}
    \item \todo{Contribution 2}
    \item \todo{Contribution 3}
\end{itemize}

%% ============================================================
\section{Related Work}
\label{sec:related}

\paragraph{<Topic 1>.}
\todo{Related work — highlight difference from our work}

\paragraph{<Topic 2>.}
\todo{Related work group 2}

%% ============================================================
\section{Background}
\label{sec:background}

% ACL papers often need linguistic/task background
\paragraph{Task Definition.}
\todo{Formal task definition}

\paragraph{Notation.}
\todo{Key notation and terminology}

%% ============================================================
\section{Method}
\label{sec:method}

\todo{Method overview}

\subsection{<Component 1>}
\todo{Architecture or algorithm details}

\subsection{<Component 2>}
\todo{Key technical contribution}

\subsection{Training}
\todo{Training objective, data augmentation if any}

%% ============================================================
\section{Experimental Setup}
\label{sec:setup}

\paragraph{Datasets.}
\todo{Dataset descriptions — include language, size, split info}

\paragraph{Baselines.}
\todo{Baseline systems}

\paragraph{Evaluation Metrics.}
\todo{Metrics with justification}

\paragraph{Implementation Details.}
\todo{Model size, hyperparameters, hardware}

%% ============================================================
\section{Results and Analysis}
\label{sec:results}

\subsection{Main Results}
\todo{Main comparison table}

\subsection{Ablation Study}
\todo{Component analysis}

\subsection{Error Analysis}
% ACL specifically values error analysis
\todo{Systematic error categorization and examples}

\subsection{Case Study}
\todo{Qualitative examples with discussion}

%% ============================================================
\section{Discussion}
\label{sec:discussion}
\todo{Broader implications, connections to linguistics}

%% ============================================================
\section{Conclusion}
\label{sec:conclusion}
\todo{Summary, limitations, future work}

%% ============================================================
% Limitations section (required by ACL since 2023)
\section*{Limitations}
\todo{Honest discussion of limitations}

% Ethics Statement (if applicable)
% \section*{Ethics Statement}
% \todo{Ethical considerations, data privacy, bias}

% Acknowledgments
% \section*{Acknowledgments}
% \todo{Funding, computing resources, helpful discussions}

\bibliography{references}
\bibliographystyle{acl_natbib}

%% ============================================================
\appendix
\section{Dataset Details}
\label{app:data}
\todo{Full dataset statistics, preprocessing}

\section{Additional Results}
\label{app:results}
\todo{Per-language/per-domain breakdowns}

\section{Annotation Guidelines}
\label{app:annotation}
\todo{If applicable — human evaluation setup}

\end{document}
